\usepackage{amsthm}
\usepackage{xparse}
\usepackage{etoolbox}
% \usepackage{tcolorbox}
% \tcbuselibrary{skins}

\makeatletter
% \patchcmd{\proof}{\@addpunct{.}}{}{}{} % <-- This patch is unreliable.
% We replace it by redefining the 'proof' environment without the dot.
\renewenvironment{proof}[1][\proofname]{%
  \par
  \normalfont \topsep6\p@\@plus6\p@\relax
  \trivlist
  \item[\hskip\labelsep
        \itshape
    #1]\ignorespaces % <--- \@addpunct{.} removed here
}{%
  \endtrivlist
}
\makeatother

% Put the current \qedsymbol as a right-aligned tag (works in equation/align/gather/split)
\newcommand{\qedtag}{\tag*{\qedsymbol}}

% For displays without tags (like \[ ... \]) fall back to a manual right push
\newcommand{\qedatdisplay}{\hfill\qedsymbol}

\newcommand{\manualqed}{\hfill\qedsymbol}

% Factory: #1=base env, #2=proof/detail env, #3=wrapper, #4=color, #5=printed title
\NewDocumentCommand{\createResEnviornment}{mmmmm}{%
  % Numbered base theorem
  \newtheorem{#1}{\sffamily\bfseries\color{#4!50!black}\MakeUppercase{#5}}[section]

  % Unnumbered base theorem
  \newtheorem*{#1*}{\sffamily\bfseries\color{#4!50!black}\MakeUppercase{#5}}

  % Proof/detail environment with optional title.
  % Proof/detail environment with optional title.
  % This avoids amsthm proof indentation.
  \NewDocumentEnvironment{#2}{ o }{%
    \par\noindent
    \IfNoValueTF{##1}{%
      \textbf{\color{#4!50!black}\LangProof}\quad
    }{%
      \IfBlankTF{##1}{%
        % Empty title: print nothing
      }{%
        \textbf{\color{#4!50!black}##1}\quad
      }%
    }%
    \renewcommand{\qedsymbol}{\textcolor{#4!40!black}{\rule{1.5ex}{1.5ex}}}%
    \ignorespaces
  }{%
    \par
  }%

  % ---------------- Numbered wrapper ----------------
  \NewDocumentEnvironment{#3}{s o}{%
    \IfBooleanTF{##1}{%
      \begin{tcolorbox}[
        enhanced,
        colback=#4!8,
        fonttitle=\bfseries,
        sharp corners,
        boxrule=0pt,
        frame hidden,
        borderline west={2pt}{0pt}{#4!70},
        breakable,
        after skip=0pt]%
    }{%
      \begin{tcolorbox}[
        enhanced,
        colback=#4!8,
        fonttitle=\bfseries,
        sharp corners,
        boxrule=0pt,
        frame hidden,
        borderline west={2pt}{0pt}{#4!70},
        breakable]%
    }%
    \IfNoValueTF{##2}{\begin{#1}}{\begin{#1}[##2]}%
    \NewDocumentCommand{\nextpart}{ o s }{%
      \end{#1}%
      \end{tcolorbox}%
      \begin{tcolorbox}[
        enhanced,
        colback=#4!2,
        fonttitle=\bfseries,
        sharp corners,
        boxrule=0pt,
        frame hidden,
        borderline west={2pt}{0pt}{#4!70},
        breakable,
        before skip=0pt]%
      \IfNoValueTF{####1}{%
        \begin{#2}%
      }{%
        \begin{#2}[####1]%
      }%
      \ignorespaces
    }%
  }{%
    \IfBooleanTF{##1}{\end{#2}}{\end{#1}}%
    \end{tcolorbox}%
  }%

  % ---------------- Unnumbered wrapper ----------------
  \NewDocumentEnvironment{#3*}{s o}{%
    \IfBooleanTF{##1}{%
      \begin{tcolorbox}[
        enhanced,
        colback=#4!8,
        fonttitle=\bfseries,
        sharp corners,
        boxrule=0pt,
        frame hidden,
        borderline west={2pt}{0pt}{#4!70},
        breakable,
        after skip=0pt]%
    }{%
      \begin{tcolorbox}[
        enhanced,
        colback=#4!8,
        fonttitle=\bfseries,
        sharp corners,
        boxrule=0pt,
        frame hidden,
        borderline west={2pt}{0pt}{#4!70},
        breakable]%
    }%
    \IfNoValueTF{##2}{\begin{#1*}}{\begin{#1*}[##2]}%
    \NewDocumentCommand{\nextpart}{ o s }{%
      \end{#1*}%
      \end{tcolorbox}%
      \begin{tcolorbox}[
        enhanced,
        colback=#4!2,
        fonttitle=\bfseries,
        sharp corners,
        boxrule=0pt,
        frame hidden,
        borderline west={2pt}{0pt}{#4!70},
        breakable,
        before skip=0pt]%
      \IfNoValueTF{####1}{%
        \begin{#2}%
      }{%
        \begin{#2}[####1]%
      }%
      \ignorespaces
    }%
  }{%
    \IfBooleanTF{##1}{\end{#2}}{\end{#1*}}%
    \end{tcolorbox}%
  }%
}

% Instantiate
\createResEnviornment{lemma}{lemproof}{nicelemma}{purple}{\LangLemma}
\createResEnviornment{theorem}{theoproof}{nicetheorem}{red}{\LangTheorem}
\createResEnviornment{corollary}{corproof}{nicecorollary}{orange}{\LangCorollary}
\createResEnviornment{definition}{defproof}{nicedefinition}{PineGreen}{\LangDefinition}
\createResEnviornment{Algorithm}{algproof}{nicealgorithm}{NavyBlue}{\LangAlgorithm}
\createResEnviornment{code}{codeproof}{nicecode}{Black}{\LangCode}
\createResEnviornment{example}{exampleproof}{niceexample}{orange}{\LangExample}
\createResEnviornment{exercise}{exerciseproof}{niceexercise}{Periwinkle}{\LangExercise}
\createResEnviornment{note}{noteproof}{nicenote}{DarkOrchid}{\LangNote}

% New property / eigenschap box
\createResEnviornment{property}{propproof}{niceproperty}{RoyalPurple}{\LangProperty}

\definecolor{factcolor}{HTML}{ef476f}
\createResEnviornment{fact}{factproof}{nicefact}{factcolor}{\LangFact}

% For referencing using Cleverref
% Format: \crefname{<env_name>}{<singular>}{<plural>}
\ifdutch
  \crefname{lemma}{lemma}{lemma's}
  \Crefname{lemma}{Lemma}{Lemma's}

  \crefname{theorem}{stelling}{stellingen}
  \Crefname{theorem}{Stelling}{Stellingen}

  \crefname{corollary}{gevolg}{gevolgen}
  \Crefname{corollary}{Gevolg}{Gevolgen}

  \crefname{definition}{definitie}{definities}
  \Crefname{definition}{Definitie}{Definities}

  \crefname{Algorithm}{algoritme}{algoritmes}
  \Crefname{Algorithm}{Algoritme}{Algoritmes}

  \crefname{code}{codeblok}{codeblokken}
  \Crefname{code}{Codeblok}{Codeblokken}

  \crefname{example}{voorbeeld}{voorbeelden}
  \Crefname{example}{Voorbeeld}{Voorbeelden}

  \crefname{exercise}{oefening}{oefeningen}
  \Crefname{exercise}{Oefening}{Oefeningen}

  \crefname{note}{opmerking}{opmerkingen}
  \Crefname{note}{Opmerking}{Opmerkingen}

  \crefname{property}{eigenschap}{eigenschappen}
  \Crefname{property}{Eigenschap}{Eigenschappen}

  \crefname{fact}{feit}{feiten}
  \Crefname{fact}{Feit}{Feiten}
\else
  \crefname{lemma}{lemma}{lemmas}
  \Crefname{lemma}{Lemma}{Lemmas}

  \crefname{theorem}{theorem}{theorems}
  \Crefname{theorem}{Theorem}{Theorems}

  \crefname{corollary}{corollary}{corollaries}
  \Crefname{corollary}{Corollary}{Corollaries}

  \crefname{definition}{definition}{definitions}
  \Crefname{definition}{Definition}{Definitions}

  \crefname{Algorithm}{algorithm}{algorithms}
  \Crefname{Algorithm}{Algorithm}{Algorithms}

  \crefname{code}{code block}{code blocks}
  \Crefname{code}{Code block}{Code blocks}

  \crefname{example}{example}{examples}
  \Crefname{example}{Example}{Examples}

  \crefname{exercise}{exercise}{exercises}
  \Crefname{exercise}{Exercise}{Exercises}

  \crefname{note}{note}{notes}
  \Crefname{note}{Note}{Notes}

  \crefname{property}{property}{properties}
  \Crefname{property}{Property}{Properties}

  \crefname{fact}{fact}{facts}
  \Crefname{fact}{Fact}{Facts}
\fi


% Stuff for code boxes
% If a dedicated JetBrains code font isn't defined in the preamble, define it here.
% (Assumes LuaLaTeX/XeLaTeX + fontspec is loaded in your preamble.)
\makeatletter
\@ifundefined{CodeFont}{%
  \newfontfamily\CodeFont{JetBrains Mono}
}{}
\makeatother

% Basic colors for syntax highlighting
\definecolor{codekeyword}{HTML}{005CC5} % blue keywords
\definecolor{codestring}{HTML}{032F62}  % dark blue strings
\definecolor{codecomment}{HTML}{6A737D} % gray comments

\lstset{
  basicstyle=\ttfamily\small,
  showstringspaces=false,
  upquote=true,
  tabsize=2,
  breaklines=true,
  columns=fullflexible,
  keepspaces=true,
  frame=none, framerule=0pt,
}

\lstdefinestyle{codebox}{
  basicstyle=\CodeFont\small,      % JetBrains only in code panel
  keywordstyle=\color{codekeyword}\bfseries,
  stringstyle=\color{codestring},
  commentstyle=\color{codecomment}\itshape,
  numbers=none,                    % no line numbers
  frame=none, framerule=0pt,       % avoid double frame
  xleftmargin=0em,
  aboveskip=4pt, belowskip=2pt,
}

% Apply JetBrains style only inside the code "proof" part
\AtBeginEnvironment{codeproof}{\lstset{style=codebox}}

% A custom colored text highlight with rounded corners

\newtcbox{\inlinebox}[1][]{
  on line,
  tcbox raise base,        % align baseline with surrounding text 
  before upper=\vphantom{AgjpqyHMU}, % includes depth for descenders
  colback=gray!15,
  colframe=gray!15,
  boxsep=0pt,
  left=2pt, right=2pt, top=0.5pt, bottom=0.25pt,
  arc=2pt,
  enhanced,
  tcbox width=auto limited,        % ← allow wrapping up to \linewidth
  halign=left,                     % nice left alignment when wrapping
  valign=center,
  #1
}

\usepackage{seqsplit}

% \inlinetext{...}[<color>]  (default gray!15)
\NewDocumentCommand{\inlinetext}{m O{gray!15}}{%
  \inlinebox[colback=#2,colframe=#2]{#1}%
}

% \inlinecode{...}[<color>]  (default gray!15)
\NewDocumentCommand{\inlinecode}{m O{gray!15}}{%
  \inlinebox[colback=#2,colframe=#2]{\texttt{#1}}%
}

% ---------- Code for Equation Boxes ------------

\NewDocumentEnvironment{equationbox}{ s O{black} O{} +b }
{%
\begingroup%
\setlength{\abovedisplayskip}{6pt}%
\setlength{\belowdisplayskip}{4pt}%
\setlength{\abovedisplayshortskip}{6pt}%
\setlength{\belowdisplayshortskip}{4pt}%
\if\relax\detokenize{#3}\relax%
    \tcbset{eqfill/.style={opacityback=0, colback=white}}%
\else%
    \tcbset{eqfill/.style={opacityback=1, colback=#3}}%
\fi%
\tcbset{eqbox/.style={%
    eqfill,%
    sharp corners,%
    boxrule=0.8pt,%
    colframe=#2,%
    top=3pt, bottom=3pt, left=5pt, right=5pt, boxsep=1pt,%
    nobeforeafter,%
    tcbox raise base,%
    enhanced%
}}%
\IfBooleanTF{#1}%
{%
    \[%
        \tcbox[eqbox, standard jigsaw]{\ensuremath{\displaystyle\ignorespaces #4\unskip}}%
    \]%
}%
{%
    \begin{equation}%
        \tcbox[eqbox, standard jigsaw]{\ensuremath{\displaystyle\ignorespaces #4\unskip}}%
    \end{equation}%
}%
\endgroup%
}
{\ignorespacesafterend}

% barbox environment: \begin{barbox}[<barcolor>][<bgcolor>] ... \end{barbox}
\NewDocumentEnvironment{barbox}{ O{gray} o }{%
  % #1 = bar color  (optional, default gray)
  % #2 = background (optional, if absent use white)
  \IfNoValueTF{#2}{%
    % No background color given -> use white (looks transparent on white paper)
    \begin{tcolorbox}[
      enhanced,
      breakable,
      colback=white,
      colframe=white,
      borderline west={3pt}{0pt}{#1},
      left=10pt,
      right=0pt,
      boxrule=0pt,
    ]%
  }{%
    % Background color given -> use it
    \begin{tcolorbox}[
      enhanced,
      breakable,
      colback=#2,
      colframe=white,
      borderline west={3pt}{0pt}{#1},
      left=10pt,
      right=0pt,
      boxrule=0pt,
    ]%
  }%
}{%
  \end{tcolorbox}%
}

% Custom round box to put text in:

% #1 = background color   (optional, default: light gray)
% #2 = border color       (optional, default: none)
\NewDocumentEnvironment{roundbox}{ O{gray!15} O{} }{%
  \IfBlankTF{#2}{%
    % no border color given
    \begin{tcolorbox}[
      enhanced,
      frame empty,
      colback=#1,
      rounded corners,
      arc=8pt,
      width=\textwidth,
      boxsep=4pt
    ]%
  }{%
    % border color given
    \begin{tcolorbox}[
      enhanced,
      colback=#1,
      colframe=#2,
      rounded corners,
      arc=8pt,
      width=\textwidth,
      boxsep=4pt
    ]%
  }%
}{%
  \end{tcolorbox}
}